\Needspace{8\baselineskip}
\section{开放模型：能力特点与运行资源}

\subsection{模型规模与架构：稠密模型和混合专家}
本次研究覆盖8项开放模型：Qwen和Gemma的4项中等规模模型，面向专业个人工作站及部门服务；DeepSeek、GLM和Kimi的4项大型模型，覆盖更复杂的研究与多步任务。理解这些模型，可以先看两个数字：需要保存多少参数，以及每次计算使用多少参数。

模型名称中的B表示十亿参数，T表示万亿参数。稠密模型在每次计算中使用其主要网络参数；混合专家模型（MoE）保存多组专家，每次选择其中一部分参与计算。例如，Qwen3.6-35B-A3B共包含约35B参数，每次激活约3B。总参数量主要影响权重存储，激活参数量主要影响每次计算，两者共同形成MoE的资源特点。

\subsection{专业个人工作站模型：Qwen与Gemma}
Qwen3.8-27B、Qwen3.6-35B-A3B、Gemma 4 31B和26B A4B都支持长文本及多模态理解。它们为财报阅读、资料问答、表格分析和工具调用提供了专业个人机级别的模型基础，也可以在部门服务器上供多人使用。

\begin{table}[H]\small\centering
\caption{中等规模模型的结构与官方权重体量}\label{tab:personal-models}
\begin{tabularx}{\linewidth}{@{}P{35mm}Y P{25mm}P{20mm}@{}}\toprule
模型 & 模型卡规模口径 & 原生上下文 & 权重体量\\\midrule
Qwen3.8-27B & 语言模型27B，稠密架构。 & 262,144 & 55.56 GB\\
Qwen3.6-35B-A3B & 总35B，每次激活3B。 & 262,144 & 71.90 GB\\
Gemma 4 31B IT & 卡片30.7B，另列视觉编码器。 & 262,144 & 62.55 GB\\
Gemma 4 26B A4B IT & 总25.2B，每次激活3.8B，另列视觉编码器。 & 262,144 & 51.61 GB\\\bottomrule
\end{tabularx}
\source{官方模型卡、配置与固定修订的权重索引。表中均为BF16制品，GB为十进制单位。\src{TECH-004}\src{TECH-013}\src{TECH-014}\src{TECH-006}}
\end{table}

\textbf{Qwen3.8-27B侧重紧凑稠密模型中的综合能力。}它支持文本、图像及视频理解，提供可调节的思考模式，并保留多轮任务的推理上下文。研究员既可以让它快速提取公告字段，也可以增加思考投入，比较材料中的证据、解释差异和组织后续操作。模型卡中的文档理解和工具任务评测，与这类工作直接相关。\src{TECH-004}

\textbf{Qwen3.6-35B-A3B侧重低激活比例的MoE路线。}它以约3B的激活参数参与每次计算，同时保有35B级模型容量，并提供多模态和工具使用能力。这种结构的价值是控制每次调用的计算量，适合持续对话、反复调用工具及多人服务；全部专家权重则约占71.90 GB的BF16存储空间。\src{TECH-013}

\textbf{Gemma 4提供Google的稠密与MoE两条路线。}31B和26B A4B均支持文本、图像、多语言及函数调用，可以连接英文财报、中文资料和图表阅读。31B侧重稠密模型的能力表现，26B A4B通过每次激活约3.8B参数控制计算量。\src{TECH-014}\src{TECH-006}

Google模型卡的128K长文检索测试（MRCR v2，8 needle，平均值）中，31B和26B A4B分别取得66.4\%和44.1\%。该任务要求从长上下文中找回指定信息，31B在这一测试中表现更好。对于同时阅读多份报告的应用，这类成绩提供了模型保持和找回材料信息的比较依据。\src{TECH-014}\src{TECH-006}

\subsection{大型模型：DeepSeek、GLM与Kimi}
大型模型扩大了研究、工具使用和长任务的能力覆盖，同时通过稀疏激活、低精度和缓存设计提高计算效率。本组模型的公开上下文配置均达到约1M，权重制品从数百GB延伸至TB级，对应高内存工作站、多GPU服务器及异构计算等运行形态。

\begin{table}[H]\small\centering
\caption{较大规模模型的结构与当前官方制品}\label{tab:shared-models}
\begin{tabularx}{\linewidth}{@{}P{34mm}Y P{24mm}P{22mm}@{}}\toprule
模型 & 规模及结构 & 制品精度 & 权重体量\\\midrule
DeepSeek-V4.1-Flash & 主干552B，另有196B条件记忆；输入处理/生成阶段激活8B/16B。 & 混合、打包 & 510.29 GB\\
GLM-5.3-Flash & 卡片320B，每次激活18B；混合注意力。 & FP8/BF16 & 328.33 GB\\
GLM-5.3 & 官方制品约753.33B；MoE，沿用GLM-5.2基模。 & FP8/BF16 & 755.62 GB\\
Kimi K3 & 总2.8T，每次激活104B；混合注意力与MoE。 & 4位等混合格式 & 1,560.86 GB\\\bottomrule
\end{tabularx}
\source{官方卡片及权重索引；四项配置均列出约1M上下文上限。权重体量按表列精度统计。\src{TECH-008}\src{TECH-016}\src{TECH-015}\src{TECH-017}}
\end{table}

\textbf{DeepSeek-V4.1-Flash的突出特点是长上下文与Agent执行效率。}它支持图像和文本输入，并将输入处理与逐步生成分为不同计算阶段，分别激活约8B和16B参数。跨层缓存共享、稀疏注意力和低精度缓存减少了保存、读取长上下文的开销，适合反复读取资料和工具结果的研究任务。其指令模型公开评测覆盖自动化、长任务和带工具的图表分析，评测采用最高思考强度。\src{TECH-008}

\textbf{GLM-5.3-Flash将多模态能力与较低计算量结合。}它采用新训练的多模态基模，每次激活约18B参数，通过线性注意力与稀疏注意力处理长材料。文本、图像和工具调用可以组成连续任务，例如读取财报页面、提取表格，再调用计算工具形成比较结果。当前权重约328.33 GB，是本组中体量较小的一项；官方同时给出了GPU和CPU/GPU异构部署路径。\src{TECH-016}

\textbf{GLM-5.3的重点是复杂工具任务与长过程执行。}它沿用GLM-5.2基模，通过后训练增强代码、自动化和多步任务能力。发布方在Toolathlon Verified中报告73.0，采用官方评测服务的三次独立运行平均单次成功率口径。对于研究系统，这类能力可以用于组织数据查询、计算脚本、文件处理及连续工具操作。\src{TECH-015}

\textbf{Kimi K3的特点是研究、办公和多模态任务覆盖较完整。}模型卡包括办公PDF、电子表格、公司金融与金融Agent评测，展示了从读取材料到操作工具、交付文件的能力。卡片引用Vals AI的CorpFin v2与Finance Agent v2成绩分别为71.6和54.4，对应最高思考强度模型列。它采用2.8T规模的MoE，每次激活104B，并以混合低精度权重控制存储；当前约1.56 TB的制品对应更大规模的集中资源配置。\src{TECH-017}

\Needspace{16\baselineskip}
\subsection{模型运行中的显存与内存需求}\label{sec:memory}
运行一个模型，内存主要用于三件事：\textbf{保存模型权重、保留当前任务的上下文、支持计算过程}。普通GPU服务器主要使用显存；Apple统一内存和CPU/GPU异构路线则让系统内存参与模型运行。图\ref{fig:memory}给出这三部分的关系。

\begin{figure}[H]\centering
\begin{tikzpicture}[node distance=5mm]
\node[flow,text width=32mm,minimum height=18mm](a){模型权重\\保存模型学到的参数};
\node[flow,right=of a,text width=32mm,minimum height=18mm](b){任务上下文\\保存已读材料的状态};
\node[flow,right=of b,text width=32mm,minimum height=18mm](c){计算空间\\保存计算中的临时数据};
\node[below=9mm of b,font=\small\bfseries\color{navy}](d){三部分共同占用显存或内存};
\draw[arr](a.south)--(d.north west);\draw[arr](b)--(d);\draw[arr](c.south)--(d.north east);
\end{tikzpicture}
\caption{模型运行的三类主要内存用途。}\label{fig:memory}
\end{figure}

\textbf{权重决定基本装载规模。}Qwen3.8-27B的BF16权重约55.56 GB，Qwen3.6-35B-A3B约71.90 GB。后者每次只选择部分专家计算，但内存中保存的是全部专家。因此，选择模型时，总体量回答“需要多大的存储空间”，激活量回答“每次调用要做多少计算”。对于48 GB显存配置，可通过降低精度、跨设备分配或CPU/GPU协作安排这两款模型。

\textbf{上下文占用随任务长度和同时运行的请求变化。}一个短问题只包含少量输入；两期财报比较则包含长篇材料、提问和工具结果。推理系统把已经处理过的信息保存为上下文状态，其中常见形式是键值缓存（KV cache），以便继续生成和追问。多人共用一份模型权重时，各个请求还会产生相应的上下文状态，共同占用剩余空间。

\textbf{计算过程也需要工作空间。}矩阵计算、图像处理和加速执行会使用临时数据与缓冲区。由此，设备配置先容纳权重，再为预期的材料长度、同时处理的请求及计算过程分配余量。

\textbf{量化通过减少数值位宽节省空间。}例如，BF16使用16位浮点表示，FP8和4位格式进一步压缩相应权重或缓存。量化让较大模型进入更紧凑的设备，也为更长上下文或更多请求留出空间；其计算效率由对应硬件和推理内核共同发挥。
